% !TEX program = xelatex
% ============================================================================
% 实验报告导言区片段（由 Lab_Workflow_Generator 提供，配方取自上游可编译样例）
%
% 用法：主文件首行写
%         \documentclass[UTF8,a4paper,10pt]{ctexart}
%       紧随其后 `% !TEX program = xelatex
% ============================================================================
% 实验报告导言区片段（由 Lab_Workflow_Generator 提供，配方取自上游可编译样例）
%
% 用法：主文件首行写
%         \documentclass[UTF8,a4paper,10pt]{ctexart}
%       紧随其后 `% !TEX program = xelatex
% ============================================================================
% 实验报告导言区片段（由 Lab_Workflow_Generator 提供，配方取自上游可编译样例）
%
% 用法：主文件首行写
%         \documentclass[UTF8,a4paper,10pt]{ctexart}
%       紧随其后 `% !TEX program = xelatex
% ============================================================================
% 实验报告导言区片段（由 Lab_Workflow_Generator 提供，配方取自上游可编译样例）
%
% 用法：主文件首行写
%         \documentclass[UTF8,a4paper,10pt]{ctexart}
%       紧随其后 `\input{report_preamble}`（或把本文件内容整段粘贴进导言区）。
%
% 依赖：XeLaTeX + ctex；宏包 booktabs/float/caption/subcaption/bicaption/
%       siunitx/tikz/fancyhdr/hyperref。
% 校徽：脚手架会把 bnu_logo.png 复制到 <实验标题>/lab_report/assets/，
%       因此默认相对路径 assets/bnu_logo.png 在 lab_report/ 目录下即有效。
% ============================================================================
\usepackage{amsmath,amssymb}
\usepackage{booktabs,longtable,tabularx,array}
\usepackage{graphicx,float}
\usepackage{geometry}
\usepackage{setspace}
\usepackage{enumitem}
\usepackage{xcolor}
\usepackage{caption}
\usepackage{subcaption}
\usepackage{bicaption}
\usepackage{tikz}
\usepackage{siunitx}
\usepackage{fancyhdr}
\usepackage[hidelinks]{hyperref}

% —— 版式：A4、四边 2cm、1.5 倍行距、首行缩进 2 字符、段间距 ——
\geometry{left=2cm,right=2cm,top=2cm,bottom=2cm}
\onehalfspacing
\setlength{\parindent}{2em}
\setlength{\parskip}{0.15em}
\setlist{nosep,leftmargin=2.2em}

% —— 物理量与不确定度（推荐写法）——
%   \SI{438.5(17)}{\mega\hertz}   → 438.5(17) MHz
%   \SI{0.35}{\percent}           → 0.35 %
%   \ang{89.8}                    → 89.8°
\sisetup{detect-all,per-mode=symbol,separate-uncertainty=true}

% —— 图题：中文主图题 + 英文 Fig. 图题（双语）——
\captionsetup{font=small,labelsep=quad}
\captionsetup[figure][bi-second]{name=Fig.,labelsep=quad}
\captionsetup[subfigure]{font=small,labelfont=bf,labelsep=space}

% —— 一级标题中文编号：一、二、三 …（与老师小论文格式一致）——
\ctexset{
  section={
    name={,、},
    number=\chinese{section},
    format=\zihao{4}\heiti\raggedright,
    beforeskip=1.1ex,
    afterskip=0.7ex
  },
  subsection={
    format=\zihao{-4}\heiti\raggedright,
    beforeskip=0.9ex,
    afterskip=0.5ex
  }
}

% —— 页眉页脚：左=实验名称，右=课程名，页脚居中页码 ——
%   生成时用本次实验名替换下面被注释的一行。
\pagestyle{fancy}
\fancyhf{}
% \fancyhead[L]{\small <实验名称>}
% \fancyhead[R]{\small 近代物理实验报告}
\fancyfoot[C]{\thepage}
\renewcommand{\headrulewidth}{0.4pt}

% —— 封面信息栏留空横线（不列班级）——
\newcommand{\blankfield}[1]{\makebox[#1]{\hrulefill}}

% —— 微积分微分号（直体 d）——
\newcommand{\dif}{\mathop{}\!\mathrm{d}}

% —— 摘要与关键词的统一样式（推荐直接调用，保证两处标题样式一致）——
\newcommand{\abstractheading}{\begin{center}{\heiti\zihao{-4} 摘\quad 要}\end{center}}
\newcommand{\keywordsline}[1]{\noindent\textbf{关键词：}#1}

% —— 首页左上角校徽：overlay 不占正文流，宽 7.0cm，
%    位于正文左页边距左移 2cm 处（页边距 2cm ⇒ 与纸张左边缘对齐），仅首页 ——
\newcommand{\bnulogofile}{assets/bnu_logo.png}
\newcommand{\bnulogo}{%
  \begin{tikzpicture}[remember picture,overlay]
    \node[anchor=north west,inner sep=0pt]
      at ([xshift=0cm,yshift=-0.90cm]current page.north west)
      {\includegraphics[width=7.0cm]{\bnulogofile}};
  \end{tikzpicture}%
}
`（或把本文件内容整段粘贴进导言区）。
%
% 依赖：XeLaTeX + ctex；宏包 booktabs/float/caption/subcaption/bicaption/
%       siunitx/tikz/fancyhdr/hyperref。
% 校徽：脚手架会把 bnu_logo.png 复制到 <实验标题>/lab_report/assets/，
%       因此默认相对路径 assets/bnu_logo.png 在 lab_report/ 目录下即有效。
% ============================================================================
\usepackage{amsmath,amssymb}
\usepackage{booktabs,longtable,tabularx,array}
\usepackage{graphicx,float}
\usepackage{geometry}
\usepackage{setspace}
\usepackage{enumitem}
\usepackage{xcolor}
\usepackage{caption}
\usepackage{subcaption}
\usepackage{bicaption}
\usepackage{tikz}
\usepackage{siunitx}
\usepackage{fancyhdr}
\usepackage[hidelinks]{hyperref}

% —— 版式：A4、四边 2cm、1.5 倍行距、首行缩进 2 字符、段间距 ——
\geometry{left=2cm,right=2cm,top=2cm,bottom=2cm}
\onehalfspacing
\setlength{\parindent}{2em}
\setlength{\parskip}{0.15em}
\setlist{nosep,leftmargin=2.2em}

% —— 物理量与不确定度（推荐写法）——
%   \SI{438.5(17)}{\mega\hertz}   → 438.5(17) MHz
%   \SI{0.35}{\percent}           → 0.35 %
%   \ang{89.8}                    → 89.8°
\sisetup{detect-all,per-mode=symbol,separate-uncertainty=true}

% —— 图题：中文主图题 + 英文 Fig. 图题（双语）——
\captionsetup{font=small,labelsep=quad}
\captionsetup[figure][bi-second]{name=Fig.,labelsep=quad}
\captionsetup[subfigure]{font=small,labelfont=bf,labelsep=space}

% —— 一级标题中文编号：一、二、三 …（与老师小论文格式一致）——
\ctexset{
  section={
    name={,、},
    number=\chinese{section},
    format=\zihao{4}\heiti\raggedright,
    beforeskip=1.1ex,
    afterskip=0.7ex
  },
  subsection={
    format=\zihao{-4}\heiti\raggedright,
    beforeskip=0.9ex,
    afterskip=0.5ex
  }
}

% —— 页眉页脚：左=实验名称，右=课程名，页脚居中页码 ——
%   生成时用本次实验名替换下面被注释的一行。
\pagestyle{fancy}
\fancyhf{}
% \fancyhead[L]{\small <实验名称>}
% \fancyhead[R]{\small 近代物理实验报告}
\fancyfoot[C]{\thepage}
\renewcommand{\headrulewidth}{0.4pt}

% —— 封面信息栏留空横线（不列班级）——
\newcommand{\blankfield}[1]{\makebox[#1]{\hrulefill}}

% —— 微积分微分号（直体 d）——
\newcommand{\dif}{\mathop{}\!\mathrm{d}}

% —— 摘要与关键词的统一样式（推荐直接调用，保证两处标题样式一致）——
\newcommand{\abstractheading}{\begin{center}{\heiti\zihao{-4} 摘\quad 要}\end{center}}
\newcommand{\keywordsline}[1]{\noindent\textbf{关键词：}#1}

% —— 首页左上角校徽：overlay 不占正文流，宽 7.0cm，
%    位于正文左页边距左移 2cm 处（页边距 2cm ⇒ 与纸张左边缘对齐），仅首页 ——
\newcommand{\bnulogofile}{assets/bnu_logo.png}
\newcommand{\bnulogo}{%
  \begin{tikzpicture}[remember picture,overlay]
    \node[anchor=north west,inner sep=0pt]
      at ([xshift=0cm,yshift=-0.90cm]current page.north west)
      {\includegraphics[width=7.0cm]{\bnulogofile}};
  \end{tikzpicture}%
}
`（或把本文件内容整段粘贴进导言区）。
%
% 依赖：XeLaTeX + ctex；宏包 booktabs/float/caption/subcaption/bicaption/
%       siunitx/tikz/fancyhdr/hyperref。
% 校徽：脚手架会把 bnu_logo.png 复制到 <实验标题>/lab_report/assets/，
%       因此默认相对路径 assets/bnu_logo.png 在 lab_report/ 目录下即有效。
% ============================================================================
\usepackage{amsmath,amssymb}
\usepackage{booktabs,longtable,tabularx,array}
\usepackage{graphicx,float}
\usepackage{geometry}
\usepackage{setspace}
\usepackage{enumitem}
\usepackage{xcolor}
\usepackage{caption}
\usepackage{subcaption}
\usepackage{bicaption}
\usepackage{tikz}
\usepackage{siunitx}
\usepackage{fancyhdr}
\usepackage[hidelinks]{hyperref}

% —— 版式：A4、四边 2cm、1.5 倍行距、首行缩进 2 字符、段间距 ——
\geometry{left=2cm,right=2cm,top=2cm,bottom=2cm}
\onehalfspacing
\setlength{\parindent}{2em}
\setlength{\parskip}{0.15em}
\setlist{nosep,leftmargin=2.2em}

% —— 物理量与不确定度（推荐写法）——
%   \SI{438.5(17)}{\mega\hertz}   → 438.5(17) MHz
%   \SI{0.35}{\percent}           → 0.35 %
%   \ang{89.8}                    → 89.8°
\sisetup{detect-all,per-mode=symbol,separate-uncertainty=true}

% —— 图题：中文主图题 + 英文 Fig. 图题（双语）——
\captionsetup{font=small,labelsep=quad}
\captionsetup[figure][bi-second]{name=Fig.,labelsep=quad}
\captionsetup[subfigure]{font=small,labelfont=bf,labelsep=space}

% —— 一级标题中文编号：一、二、三 …（与老师小论文格式一致）——
\ctexset{
  section={
    name={,、},
    number=\chinese{section},
    format=\zihao{4}\heiti\raggedright,
    beforeskip=1.1ex,
    afterskip=0.7ex
  },
  subsection={
    format=\zihao{-4}\heiti\raggedright,
    beforeskip=0.9ex,
    afterskip=0.5ex
  }
}

% —— 页眉页脚：左=实验名称，右=课程名，页脚居中页码 ——
%   生成时用本次实验名替换下面被注释的一行。
\pagestyle{fancy}
\fancyhf{}
% \fancyhead[L]{\small <实验名称>}
% \fancyhead[R]{\small 近代物理实验报告}
\fancyfoot[C]{\thepage}
\renewcommand{\headrulewidth}{0.4pt}

% —— 封面信息栏留空横线（不列班级）——
\newcommand{\blankfield}[1]{\makebox[#1]{\hrulefill}}

% —— 微积分微分号（直体 d）——
\newcommand{\dif}{\mathop{}\!\mathrm{d}}

% —— 摘要与关键词的统一样式（推荐直接调用，保证两处标题样式一致）——
\newcommand{\abstractheading}{\begin{center}{\heiti\zihao{-4} 摘\quad 要}\end{center}}
\newcommand{\keywordsline}[1]{\noindent\textbf{关键词：}#1}

% —— 首页左上角校徽：overlay 不占正文流，宽 7.0cm，
%    位于正文左页边距左移 2cm 处（页边距 2cm ⇒ 与纸张左边缘对齐），仅首页 ——
\newcommand{\bnulogofile}{assets/bnu_logo.png}
\newcommand{\bnulogo}{%
  \begin{tikzpicture}[remember picture,overlay]
    \node[anchor=north west,inner sep=0pt]
      at ([xshift=0cm,yshift=-0.90cm]current page.north west)
      {\includegraphics[width=7.0cm]{\bnulogofile}};
  \end{tikzpicture}%
}
`（或把本文件内容整段粘贴进导言区）。
%
% 依赖：XeLaTeX + ctex；宏包 booktabs/float/caption/subcaption/bicaption/
%       siunitx/tikz/fancyhdr/hyperref。
% 校徽：脚手架会把 bnu_logo.png 复制到 <实验标题>/lab_report/assets/，
%       因此默认相对路径 assets/bnu_logo.png 在 lab_report/ 目录下即有效。
% ============================================================================
\usepackage{amsmath,amssymb}
\usepackage{booktabs,longtable,tabularx,array}
\usepackage{graphicx,float}
\usepackage{geometry}
\usepackage{setspace}
\usepackage{enumitem}
\usepackage{xcolor}
\usepackage{caption}
\usepackage{subcaption}
\usepackage{bicaption}
\usepackage{tikz}
\usepackage{siunitx}
\usepackage{fancyhdr}
\usepackage[hidelinks]{hyperref}

% —— 版式：A4、四边 2cm、1.5 倍行距、首行缩进 2 字符、段间距 ——
\geometry{left=2cm,right=2cm,top=2cm,bottom=2cm}
\onehalfspacing
\setlength{\parindent}{2em}
\setlength{\parskip}{0.15em}
\setlist{nosep,leftmargin=2.2em}

% —— 物理量与不确定度（推荐写法）——
%   \SI{438.5(17)}{\mega\hertz}   → 438.5(17) MHz
%   \SI{0.35}{\percent}           → 0.35 %
%   \ang{89.8}                    → 89.8°
\sisetup{detect-all,per-mode=symbol,separate-uncertainty=true}

% —— 图题：中文主图题 + 英文 Fig. 图题（双语）——
\captionsetup{font=small,labelsep=quad}
\captionsetup[figure][bi-second]{name=Fig.,labelsep=quad}
\captionsetup[subfigure]{font=small,labelfont=bf,labelsep=space}

% —— 一级标题中文编号：一、二、三 …（与老师小论文格式一致）——
\ctexset{
  section={
    name={,、},
    number=\chinese{section},
    format=\zihao{4}\heiti\raggedright,
    beforeskip=1.1ex,
    afterskip=0.7ex
  },
  subsection={
    format=\zihao{-4}\heiti\raggedright,
    beforeskip=0.9ex,
    afterskip=0.5ex
  }
}

% —— 页眉页脚：左=实验名称，右=课程名，页脚居中页码 ——
%   生成时用本次实验名替换下面被注释的一行。
\pagestyle{fancy}
\fancyhf{}
% \fancyhead[L]{\small <实验名称>}
% \fancyhead[R]{\small 近代物理实验报告}
\fancyfoot[C]{\thepage}
\renewcommand{\headrulewidth}{0.4pt}

% —— 封面信息栏留空横线（不列班级）——
\newcommand{\blankfield}[1]{\makebox[#1]{\hrulefill}}

% —— 微积分微分号（直体 d）——
\newcommand{\dif}{\mathop{}\!\mathrm{d}}

% —— 摘要与关键词的统一样式（推荐直接调用，保证两处标题样式一致）——
\newcommand{\abstractheading}{\begin{center}{\heiti\zihao{-4} 摘\quad 要}\end{center}}
\newcommand{\keywordsline}[1]{\noindent\textbf{关键词：}#1}

% —— 首页左上角校徽：overlay 不占正文流，宽 7.0cm，
%    位于正文左页边距左移 2cm 处（页边距 2cm ⇒ 与纸张左边缘对齐），仅首页 ——
\newcommand{\bnulogofile}{assets/bnu_logo.png}
\newcommand{\bnulogo}{%
  \begin{tikzpicture}[remember picture,overlay]
    \node[anchor=north west,inner sep=0pt]
      at ([xshift=0cm,yshift=-0.90cm]current page.north west)
      {\includegraphics[width=7.0cm]{\bnulogofile}};
  \end{tikzpicture}%
}
